\chapter{Le constructeur de carnet d'ordres}\label{ch:ll:the-order-book-builder}

Sur un instrument à petit tick, un carnet tenu dans des arbres équilibrés demande environ \qty{70}{\nano\second} par mise à jour sur
l'ordinateur portable de ce livre~; un tableau de niveaux de prix autour des meilleurs cours en demande environ \qty{30}{\nano\second}. Le tableau l'emporte, jusqu'au jour où
l'action varie de trente pour cent~: un tableau de 256 ticks doit alors être déplacé et rempli à nouveau presque vingt mille fois, et le
carnet de la journée coûte \qty{0.7}{\second} de temps processeur au lieu de quatre millisecondes. Le \omterm{def:ll:the-order-book-builder:builder}{constructeur de carnet} transforme les événements du \omterm{def:ll:the-feed-handler:handler}{feed
handler} en la réponse que chaque stratégie demande en premier, quel est le meilleur prix et quelle quantité s'y trouve, et la réponse
doit être juste, rapide, et rester rapide le jour qui compte. Ce chapitre précise à quoi le carnet doit répondre, compare les
structures qui y répondent selon le type de flux et le régime de tick, construit un carnet ordre par ordre sur une table des ordres à \omterm{def:ll:the-order-book-builder:open}{adressage ouvert} et une
échelle de prix à recentrage, et le teste face à deux \omterm{def:ll:the-order-book-builder:builder}{constructeurs de carnet} écrits plus tôt dans la série.

\section{Ce que le carnet doit fournir, et à quelle vitesse}

\begin{definition}[Constructeur de carnet]\label{def:ll:the-order-book-builder:builder}
Un \emph{constructeur de carnet}\index{constructeur de carnet} (book builder) maintient le carnet d'ordres à cours limité (livre~10, chapitre~1) d'un instrument à partir des
événements du \omterm{def:ll:the-feed-handler:handler}{feed handler}~: avec un flux market-by-order (livre~1, chapitre~19), il conserve chaque ordre au repos et les
agrège en niveaux de prix~; avec un flux market-by-price, il conserve directement les niveaux. Il répond aux requêtes de niveau 1 et de niveau 2
(livre~1, chapitre~28) après chaque événement, et sait quand son carnet ne mérite plus confiance.
\end{definition}

Une stratégie pose trois types de question. Niveau~1~: les meilleurs cours acheteur et vendeur, avec leurs tailles, après chaque événement. Le microprix du
livre~7, chapitre~8, n'a besoin de rien d'autre. Niveau~2, les premiers niveaux de chaque côté, pour le déséquilibre de profondeur et les estimations de
la position dans la file. Et, pour ses propres ordres, niveau~3~: où se trouve mon ordre dans sa file. Le constructeur sert les trois à partir d'une seule structure,
au rythme d'une mise à jour par événement du flux, et la mise à jour est sur le chemin critique du budget du chapitre~1~: un carnet qui
met \qty{100}{\nano\second} de plus par événement retarde chaque décision d'autant, et à l'ouverture, quand les événements arrivent
par rafales (chapitre~18), il les retarde de la file d'événements qui les précède.

Chaque événement désigne un ordre par sa référence, si bien que chaque mise à jour commence par une recherche~: quel ordre, à quel prix et de quel côté.
Il modifie ensuite un niveau de prix~: il y ajoute une quantité, la réduit, ou le retire lorsqu'il se vide. Un ajout à un nouveau prix crée un niveau~;
un remplacement (spécification de Nasdaq~: l'ordre d'origine «~les actions restantes \dots\ doivent être retirées~» et une nouvelle référence
«~sera utilisée désormais~») retire un ordre et en ajoute un autre. Les deux coûts d'un carnet sont donc la recherche de l'ordre et la
mise à jour du niveau, et les structures ci-dessous diffèrent sur ces deux points.

\section{Structures selon le type de flux et le régime de tick}

\begin{definition}[Échelle de prix dense]\label{def:ll:the-order-book-builder:ladder}
Une \emph{échelle de prix dense} conserve les niveaux de prix d'un carnet dans un tableau indexé par la distance en ticks à un prix
de référence, une case par tick, qu'un niveau y soit présent ou non, couvrant une fenêtre autour des meilleurs cours~; les niveaux hors de la fenêtre
vivent dans une structure éparse, et la fenêtre est déplacée (recentrée) lorsque les meilleurs cours approchent de son bord.
\end{definition}

\begin{omfigure}
\begin{tikzpicture}[font=\scriptsize,>=Stealth]
\def\u{0.38}
\foreach \i in {0,...,23} {
  \pgfmathtruncatemacro{\q}{(\i==6)||(\i==8)||(\i==9)||(\i==11)}
  \pgfmathtruncatemacro{\a}{(\i==13)||(\i==14)||(\i==16)||(\i==19)}
  \ifnum\q=1 \draw[fill=omDef!40] (\u*\i,0) rectangle ++(\u,0.5);
  \else\ifnum\a=1 \draw[fill=omThm!35] (\u*\i,0) rectangle ++(\u,0.5);
  \else \draw[fill=gray!6,draw=gray!50] (\u*\i,0) rectangle ++(\u,0.5); \fi\fi
}
\fill[gray!25,opacity=0.5] (0,0) rectangle (\u*6,0.5);
\fill[gray!25,opacity=0.5] (\u*18,0) rectangle (\u*24,0.5);
\draw[decorate,decoration={brace,amplitude=4pt},thick] (\u*6,0.6) -- (\u*18,0.6) node[midway,above=5pt]{moitié centrale~: les meilleurs cours restent ici};
\draw[->,omDef,thick] (\u*11.5,-0.7) node[below left=-2pt]{meilleur bid (indice 11)} -- (\u*11.5,-0.05);
\draw[->,omThm,thick] (\u*13.5,-0.7) node[below right=-2pt]{meilleur ask (indice 13)} -- (\u*13.5,-0.05);
\node[draw,rounded corners=2pt,fill=gray!8,align=center,anchor=west] (far) at (\u*24+0.4,0.25) {table éparse~:\\ niveaux éloignés};
\node[draw,rounded corners=2pt,fill=omMeth!15,align=center,anchor=north west] (map) at (0,-1.6) {table des ordres (adressage ouvert)~:\\ référence $\to$ emplacement dans le pool};
\node[draw,rounded corners=2pt,fill=omMeth!8,align=center,anchor=west] (pool) at ([xshift=0.8cm]map.east) {pool d'ordres~:\\ prix, quantité, côté};
\draw[->] (map) -- (pool);
\node[anchor=north west,align=left] at (0,-2.7) {prix $=$ base $+$ indice $\times$ tick};
\end{tikzpicture}

\omcaption{Le carnet du constructeur. Les niveaux vivent dans un tableau dense indexé par les ticks à partir d'un prix de référence (bleu~: bids, rouge~: asks, gris~:
vides)~; les meilleurs bid et ask sont des indices tenus à jour. Lorsqu'un meilleur prix quitte la moitié centrale du tableau, la base est
déplacée pour recentrer les meilleurs cours, et les ordres sont de nouveau rangés. Les ordres sont retrouvés par une table à \omterm{def:ll:the-order-book-builder:open}{adressage ouvert} qui relie leur
référence à un emplacement dans un pool de taille fixe.}\label{fig:ll:the-order-book-builder:ladder}
\end{omfigure}

Trois structures se disputent les niveaux. Un arbre équilibré par côté, la réponse classique des manuels (\texttt{std::map}, dont la recherche,
la suppression et l'insertion sont logarithmiques et qui est «~généralement implémenté sous forme d'arbres rouge-noir~»), alloue un nœud pour chaque nouveau
niveau et en parcourt quelques-uns à chaque mise à jour. Un vecteur de niveaux, trié du meilleur au moins bon, trouve un niveau par recherche dichotomique et garde
les meilleurs cours en tête, mais insérer un niveau au milieu décale tous les niveaux qui suivent. Une échelle dense
(la \cref{fig:ll:the-order-book-builder:ladder}) trouve un niveau par une soustraction et une division, sans jamais allouer de mémoire~; ses
coûts se situent ailleurs~: la mémoire de toute la fenêtre, le balayage jusqu'au niveau non vide suivant lorsque le meilleur se vide, et le
recentrage lorsque les meilleurs cours dérivent.

Le type de flux et le régime de tick décident. Sur un instrument à grand tick, dont la fourchette vaut un tick la plupart du temps et dont
les niveaux sont tous occupés près des meilleurs cours, toutes les structures sont rapides, car le carnet est petit et les meilleurs cours bougent rarement~: sur la
minute du simulateur du chapitre~18, les trois constructeurs mettent environ \qty{28}{\nano\second} (échelle), \qty{36}{\nano\second}
(vecteur trié) et \qty{37}{\nano\second} (arbre) à la médiane. Sur un instrument à petit tick, dont les ordres s'étalent sur des centaines
de ticks avec des trous entre eux, l'arbre et le vecteur paient le coût d'un carnet plus grand, de 65 à \qty{72}{\nano\second}, tandis que
l'échelle reste autour de \qty{31}{\nano\second} (la \cref{fig:ll:the-order-book-builder:update}). La queue de l'échelle, environ
\qty{150}{\nano\second} au p99 sur le flux à grand tick, est son balayage~: lorsque le meilleur niveau se vide, le suivant peut être
à plusieurs cases vides de distance.

\begin{omfigure}
\begin{tikzpicture}
\begin{axis}[width=12cm,height=5.6cm,ybar,bar width=9pt,ymin=0,ymax=320,
  symbolic x coords={large-tree,large-vector,large-ladder,small-tree,small-vector,small-ladder},xtick=data,
  xticklabels={{arbre},{vecteur\\ trié},{échelle},{arbre},{vecteur\\ trié},{échelle}},
  xticklabel style={align=center,font=\footnotesize},enlarge x limits=0.1,ylabel={ns par événement},
  tick label style={font=\small},label style={font=\small},grid=major,grid style={gray!25},
  legend columns=2,legend style={font=\footnotesize,at={(0.5,-0.3)},anchor=north,draw=none,fill=none,
  /tikz/every even column/.append style={column sep=8pt}},table/col sep=comma]
\addplot[fill=omDef!70,draw=omDef] table[x=key,y=q50]{figdata/low-latency/19-the-order-book-builder/measured_update_keyed.csv};
\addplot[fill=omThm!45,draw=omThm] table[x=key,y=q99]{figdata/low-latency/19-the-order-book-builder/measured_update_keyed.csv};
\legend{médiane,p99}
\node[font=\scriptsize,anchor=north] at (axis cs:large-vector,315) {grand tick (simulateur)};
\node[font=\scriptsize,anchor=north] at (axis cs:small-vector,315) {petit tick (synthétique)};
\end{axis}
\end{tikzpicture}

\omcaption{Temps d'application d'un événement avec trois structures de carnet, sur la minute du simulateur du chapitre~18 (un
instrument à grand tick, environ 190\,000 événements) et sur 100\,000 événements synthétiques à petit tick (ordres jusqu'à 300 ticks du mid). Le
coût propre du chronomètre est soustrait. Mesuré sur un ordinateur portable (Intel Core Ultra 7 155H) sous WSL2, sans cœurs isolés. Données~:
\texttt{bench\_books.py}.}\label{fig:ll:the-order-book-builder:update}
\end{omfigure}

\section{Recherche d'ordres~: tables de hachage et adressage ouvert}

\begin{definition}[Adressage ouvert]\label{def:ll:the-order-book-builder:open}
L'\emph{adressage ouvert}\index{adressage ouvert} range les entrées d'une table de hachage dans un seul tableau plat~: une clé va à l'emplacement que désigne son hachage
ou, si cet emplacement est déjà pris, au premier emplacement libre d'une séquence de sondage fixée (sondage linéaire~: les emplacements qui suivent), et une
recherche suit la même séquence jusqu'à trouver la clé ou un emplacement vide.
\end{definition}

\texttt{std::unordered\_map}, la table de hachage de la bibliothèque standard, conserve chaque entrée dans un nœud qui lui est propre, alloué à l'insertion
et atteint par le pointeur d'un seau~: une allocation par ordre et un pointeur à suivre à chaque recherche. L'\omterm{def:ll:the-order-book-builder:open}{adressage ouvert} conserve
les références aux ordres et leurs emplacements dans le pool dans deux tableaux plats, dimensionnés une fois pour toutes au démarrage à au moins le double du nombre maximal
d'ordres vivants attendus~; une recherche se réduit à une multiplication, un masque et, en général, une \omterm{def:ll:memory-hierarchy-and-caches:line}{ligne de cache}. La suppression est la partie délicate~: vider
un emplacement couperait les séquences de sondage qui le traversent~; la construction décale donc vers l'arrière les entrées suivantes de la suite dans le trou
lorsque leur emplacement d'origine le permet, au lieu de laisser une pierre tombale, et les recherches restent aussi brèves après une journée d'annulations que
à l'ouverture.

\omcode{code/firm/bookbuilder/cpp/firm_bookbuilder.hpp}{37}{53}{Suppression par décalage arrière~: les séquences de sondage restent intactes
sans pierres tombales.}\label{lst:ll:the-order-book-builder:erase}

Les ordres eux-mêmes résident dans un pool d'enregistrements de taille fixe alloués au démarrage, avec une \omterm{def:ll:cpp-for-latency-i-memory:pool}{liste libre} (\omterm{def:ll:cpp-for-latency-i-memory:pool}{liste libre} (free list), chapitre~6)~; la table stocke
un indice d'emplacement, non un pointeur. Après le démarrage, appliquer un événement n'alloue rien~: le test de la construction compte les allocations du tas
pendant que les deux jeux de test sont appliqués, et n'en trouve aucune.

\section{Exactitude en présence de trous, de carnets croisés et de reprise}

\begin{definition}[Invariant de carnet]\label{def:ll:the-order-book-builder:invariant}
Un \emph{invariant de carnet} est une propriété que tout carnet correct vérifie après chaque événement, contrôlée par le constructeur dans ses versions de débogage~:
la quantité de chaque niveau égale la somme des quantités de ses ordres, aucun niveau vide n'est signalé, et le meilleur cours acheteur est inférieur
au meilleur cours vendeur.
\end{definition}

Un carnet est bien plus souvent faux à cause de ses entrées que de sa structure. Un événement peut désigner un ordre inconnu (une
référence perdue dans un trou), réduire un ordre de plus qu'il ne contient, ou arriver alors que le \omterm{def:ll:the-feed-handler:handler}{feed handler} a marqué le carnet comme périmé.
Les règles de la construction suivent le \omterm{def:ll:the-feed-handler:handler}{feed handler} du chapitre~18~: pendant un instantané (entre ses événements de début et de fin), le carnet est
vidé, reconstruit à partir des ordres de l'instantané et marqué comme périmé~; un événement pour un ordre inconnu est ignoré et compté~; une
réduction supérieure à l'ordre retire l'ordre. Les invariants attrapent le reste~: la référence Python les vérifie après chaque
événement des deux jeux de test, et la méthode \texttt{check()} du constructeur C++ recalcule chaque niveau à partir du pool.

Un carnet croisé mérite un mot. La plateforme n'en publie jamais en négociation continue, donc un carnet croisé dans le constructeur signifie
un événement perdu ou un bogue, et doit être traité comme périmé, sans négocier dessus. Lors des enchères et dans certaines vues consolidées, les croisements sont
légitimes~; le constructeur les signale alors au lieu de les corriger.

L'échelle présente un seul risque d'exactitude qui lui est propre, le recentrage. Lorsque le meilleur cours quitte la moitié centrale de la fenêtre, le constructeur
déplace la base pour que le mid soit au centre et reclasse chaque ordre vivant~: le seul parcours de tous les ordres, effectué
sur place.

\omcode{code/firm/bookbuilder/cpp/firm_bookbuilder.hpp}{220}{244}{Recentrage~: lorsqu'un meilleur prix quitte la moitié centrale, centrer
l'échelle sur le mid et reclasser les ordres vivants.}\label{lst:ll:the-order-book-builder:recentre}

À quelle fréquence~? Pour un prix qui évolue comme un mouvement brownien, un seul temps de sortie moyen suffit à y répondre.

\begin{proposition}[Recentrages d'une échelle]\label{prop:ll:the-order-book-builder:recentre}
Soit le prix, en ticks, un mouvement brownien d'écart type $\sigma$ par jour, et soit l'échelle recentrée
chaque fois que le prix s'écarte de $a$ ticks du centre ($a$ est le quart de la largeur). Le temps entre deux recentrages a pour moyenne
$a^2/\sigma^2$ jours, donc le taux à long terme est de $\sigma^2/a^2$ recentrages par jour. Un prix qui suit une tendance de $D$ ticks en une journée provoque
environ $D/a$ recentrages.
\end{proposition}

\begin{proof}
Soit $X_t$ le déplacement depuis le centre et $\tau$ le premier instant où $|X_\tau| = a$. Comme $X_t^2 - \sigma^2 t$ est une martingale et
$\tau$ admet une moyenne finie, le théorème d'arrêt optionnel donne $\E[X_\tau^2] = \sigma^2 \E[\tau]$, et $X_\tau^2 = a^2$, donc $\E[\tau] =
a^2/\sigma^2$. Les recentrages forment un processus de renouvellement dont le taux à long terme est $1/\E[\tau]$. Une tendance régulière parcourt $a$ ticks
entre deux recentrages, soit $D/a$ d'entre eux pour $D$ ticks.
\end{proof}

Une action à 100 avec un tick de $10^{-4}$ qui évolue de 1~\% par jour a $\sigma = 10\,000$ ticks~: une échelle de 4\,096 ticks
($a = 1\,024$) recentre environ 95 fois par jour, une de 65\,536 ticks environ une fois tous les trois jours. La même action qui évolue de trente
pour cent en une journée parcourt 300\,000 ticks~: environ 290 recentrages avec l'échelle la plus étroite, et chacun d'eux reclasse chaque
ordre vivant. La \Cref{fig:ll:the-order-book-builder:width}
mesure ce jour-là, sur des événements synthétiques dont le mid dérive de trente pour cent~: 256 recentrages avec 4\,096 ticks, quinze avec
65\,536, aucun avec un million, et un temps de traitement total qui tombe de \qty{0.7}{\second} avec une échelle de 256 ticks à
\qty{4}{\milli\second} pour 65\,536 ticks, puis remonte à mesure que l'échelle dépasse les caches.

\begin{omfigure}
\begin{tikzpicture}
\begin{axis}[width=11.5cm,height=5.6cm,xmode=log,ymode=log,log basis x=2,xmin=180,xmax=12000000,ymin=1,ymax=3000,
  xlabel={largeur de l'échelle (ticks, échelle logarithmique)},ylabel={ms pour la journée (échelle logarithmique)},
  xtick={256,1024,4096,16384,65536,262144,1048576,4194304},
  xticklabels={256,1\,024,4\,096,16\,384,65\,536,$2^{18}$,$2^{20}$,$2^{22}$},
  tick label style={font=\small},label style={font=\small},grid=major,grid style={gray!25},
  nodes near coords,point meta=explicit symbolic,every node near coord/.append style={font=\scriptsize,anchor=south west},
  table/col sep=comma]
\addplot[omDef,very thick,mark=*,mark size=1.6pt] table[x=width,y=total_ms,meta=recentrings]{figdata/low-latency/19-the-order-book-builder/measured_width.csv};
\end{axis}
\end{tikzpicture}

\omcaption{La journée où une action à petit tick évolue de trente pour cent (100\,000 événements synthétiques, le mid dérivant de 300\,000 ticks)~:
temps de traitement du carnet en échelle par largeur, chaque point étant annoté de son nombre de recentrages. Mesuré sur un ordinateur portable (Intel
Core Ultra 7 155H) sous WSL2, sans cœurs isolés. Données~: \texttt{bench\_books.py}.}\label{fig:ll:the-order-book-builder:width}
\end{omfigure}

\section{Tester un carnet}

Un \omterm{def:ll:the-order-book-builder:builder}{constructeur de carnet} se teste contre d'autres carnets. La référence Python du constructeur, un dictionnaire d'ordres et un de niveaux par côté
sans rien d'astucieux, est comparée événement par événement au carnet du livre~1, chapitre~28 (\texttt{firm.feed}),
sur son échantillon de messages ordre par ordre~: les mêmes dix niveaux de chaque côté après chaque message. Elle est comparée au carnet
qui alimente le moteur de variables du livre~7 (\texttt{firm.lobfeat}) sur le jeu de test de ce moteur~: les mêmes meilleurs prix et tailles après
chaque message. Les échelles C++ et Rust sont ensuite comparées à la référence Python sur deux jeux de test, les événements du simulateur
et un flux synthétique à petit tick, par un hachage 64 bits des cinq premiers niveaux de chaque côté, enchaîné sur chaque événement~; la
comparaison est menée avec l'échelle ordinaire et avec une échelle étroite qui recentre et déverse des ordres dans sa table creuse, car
c'est dans les chemins rares d'une structure que vivent ses bogues. Et les constructions à arbre et à vecteur trié du banc d'essai sont vérifiées
contre l'échelle de la même façon avant que l'une d'elles soit chronométrée.

\section{Tutoriel~: trois carnets, deux régimes}\label{tut:ll:the-order-book-builder:tut}

\begin{tutorial}
\textbf{Objectif.} Construire le carnet en échelle, le valider face aux carnets précédents, et le chronométrer face à un arbre et à un vecteur trié sur un
instrument à grand tick et un instrument à petit tick. \textbf{État final~:}
les \cref{fig:ll:the-order-book-builder:update,fig:ll:the-order-book-builder:width}, et des tests au vert en Python, C++ et Rust.
\begin{tutsteps}
\item \textbf{Les jeux de test.} \texttt{make\_book\_fixtures.py} écrit les événements du \omterm{def:ll:the-feed-handler:handler}{feed handler} pour les trois secondes du simulateur
(chapitre~18) et 6\,000 événements synthétiques à petit tick, ainsi que les hachages de niveau 2 que la référence calcule sur ceux-ci.
\item \textbf{Tests différentiels.} La référence face à \texttt{firm.feed} du livre~1 et à \texttt{firm.lobfeat} du livre~7~: les échelles C++
et Rust face à la référence (hachage après chaque événement), avec des largeurs de 4\,096 et 256 ticks~; la suppression dans la table des ordres
sur des clés qui entrent en collision~; les invariants après chaque événement~; zéro allocation après l'instanciation.
\item \textbf{Mesurer} avec \texttt{python bench\_books.py}~: les trois structures sur la minute du simulateur et sur 100\,000
événements à petit tick, puis l'échelle par largeur sur la journée où le mid dérive de trente pour cent.
\end{tutsteps}
\textbf{Pour aller plus loin.} Remplacer le parcours de l'échelle à la recherche du prochain niveau non vide par une bitmap des niveaux occupés et une
instruction count-trailing-zeros, et mesurer à nouveau le p99 sur le flux à grand tick.
\end{tutorial}

\section{Construction~: le constructeur de carnet}\label{bld:ll:the-order-book-builder:bookbuilder}

\begin{build}
\bfield{Rôle} Le carnet qui sous-tend chaque stratégie du système de trading~: il consomme les événements du \omterm{def:ll:the-feed-handler:handler}{feed handler} (chapitre~18) et
sert le niveau 1 et le niveau 2 au moteur de stratégies (chapitre~20) ainsi qu'à la barrière de risque (chapitre~22).
\bfield{Interface} Référence Python \texttt{firm\_bookbuilder.Book}~: \texttt{apply(event)}, \texttt{best()},
\texttt{depth(side, n)}, \texttt{stale}, \texttt{check()}~; \texttt{l2\_hash}~; \texttt{recentrings},
\texttt{expected\_recentrings}. C++20 \texttt{firm::book::LadderBook(tick, width, max\_orders)}~: \texttt{apply},
\texttt{best}, \texttt{depth}, \texttt{check}, \texttt{stale}, \texttt{recentrings}~; \texttt{OrderMap}~; \texttt{l2\_hash}. Rust
\texttt{firm\_bookbuilder::LadderBook} et \texttt{OrderMap}, avec les mêmes.
\bfield{Règles} Recherche d'ordres par \omterm{def:ll:the-order-book-builder:open}{adressage ouvert} dans un pool de taille fixe~; niveaux dans une échelle dense avec un repli épars~;
échelle recentrée lorsqu'un meilleur prix quitte sa moitié centrale~; carnet vidé et marqué périmé du début à la fin
d'un instantané~; aucune allocation par événement après la construction, sauf si un ordre tombe hors de l'échelle.
\bfield{Tests d'acceptation} \texttt{code/firm/bookbuilder/}~: niveau 2 identique au \texttt{firm.feed} du livre~1 après chaque message
de son échantillon et au \texttt{firm.lobfeat} du livre~7 après chaque message de son jeu de test (référence)~; les échelles C++ et Rust
égales à la référence après chaque événement des deux jeux, avec une échelle normale et une échelle étroite~; invariants après chaque événement~; zéro
allocation~; la suppression par décalage arrière de la table des ordres~; le modèle de recentrage face à un prix suivant un mouvement brownien simulé.
\bfield{Pour aller plus loin} Une bitmap des niveaux occupés pour le parcours~; le niveau~3 avec les positions dans la file pour les propres ordres de la société~;
un mode market-by-price~; un recentrage en arrière-plan dans une seconde échelle.
\end{build}

\begin{omsources}
\item Nasdaq, spécification de \emph{TotalView-ITCH 5.0}, section 1.4.5 (Order Replace).
\item cppreference.com, \texttt{std::map}.
\item Les carnets d'ordres des livres précédents~: \texttt{firm.feed} (livre~1, chapitre~28) et \texttt{firm.lobfeat} (livre~7,
chapitre~8).
\end{omsources}

\section{Exercices}

\begin{exercise}[$\star$]\label{exo:ll:the-order-book-builder:1}
Une échelle de 4\,096 ticks avec un tick de 0.01 est centrée sur 50.00. Quel est son prix de base, quel indice contient 49.73, et où
doit se trouver le meilleur cours acheteur pour que l'échelle reste en place~?
\end{exercise}

\begin{exercise}[$\star$]\label{exo:ll:the-order-book-builder:2}
Quelle mémoire occupe une échelle de 65\,536 ticks avec 16 octets par niveau et deux côtés~? Tient-elle dans le cache de deuxième niveau
du chapitre~3~?
\end{exercise}

\begin{exercise}[$\star$]\label{exo:ll:the-order-book-builder:3}
Un ordre de 500 à 20.10 est remplacé par un ordre de 300 à 20.11. Quels événements le carnet applique-t-il, et que devient
la place de l'ordre dans sa file~?
\end{exercise}

\begin{exercise}[$\star\star$]\label{exo:ll:the-order-book-builder:4}
Les clés 1\,024, 2\,048 et 3\,072 ont le même emplacement initial, la case 5, dans une table de 16. Elles sont insérées dans cet ordre, puis
1\,024 est effacée. Montrer la table avant et après le décalage arrière. Qu'aurait coûté une pierre tombale à la place~?
\end{exercise}

\begin{exercise}[$\star\star$]\label{exo:ll:the-order-book-builder:5}
Un titre se négocie à 40 avec un tick de 0.01 et une volatilité quotidienne de 2~\%. Avec une échelle de 1\,024 ticks, combien de
recentrages par jour prévoit la \cref{prop:ll:the-order-book-builder:recentre}~? Et avec 4\,096~?
\end{exercise}

\begin{exercise}[$\star\star$]\label{exo:ll:the-order-book-builder:6}
Le constructeur reçoit une réduction de 300 pour un ordre qui en détient 200. Que fait-il, et que signale-t-il à la surveillance~?
\end{exercise}

\begin{exercise}[$\star\star\star$]\label{exo:ll:the-order-book-builder:7}
\emph{Programmation.} Ajouter une bitmap des niveaux occupés à l'échelle C++ (un bit par tick et par côté) et utiliser count-trailing-zeros pour
trouver le prochain meilleur niveau~; vérifier les hachages des jeux de test et mesurer le p99 à grand tick.
\end{exercise}

\begin{exercise}[$\star\star\star$]\label{exo:ll:the-order-book-builder:8}
\emph{Trouver la faille.} «~Nous avons dimensionné l'échelle à 512 ticks parce qu'elle tient dans le cache de premier niveau, et nos tests
un jour calme n'ont montré aucun recentrage.~»
\end{exercise}

\section{Problème~: l'action qui a bougé de trente pour cent}

\begin{problem}[{Problème du week-end --- quelle largeur pour une échelle}]\label{pb:ll:the-order-book-builder:1}
Une action à petit tick se négocie à 100 avec un tick de $10^{-4}$. Un jour ordinaire, sa volatilité est de 1.5~\%~; le jour de ce
problème, elle bouge de trente pour cent. Utiliser la \cref{prop:ll:the-order-book-builder:recentre} et les mesures de la
\cref{fig:ll:the-order-book-builder:width} (\texttt{measured\_width.csv}).

\medskip\noindent\textbf{Partie I --- Le jour ordinaire.}
\begin{enumerate}
\item Quelle est la valeur de $\sigma$ en ticks par jour~?
\item Combien de recentrages par jour une échelle de 4\,096 ticks prévoit-elle~? Et une de 65\,536~?
\item Quelle largeur correspond à un recentrage par jour~?
\item Pourquoi le jour ordinaire est-il un mauvais guide pour le dimensionnement~?
\end{enumerate}

\medskip\noindent\textbf{Partie II --- Le jour qui compte.}
\begin{enumerate}[resume]
\item Combien de ticks le prix parcourt-il, et combien de recentrages une tendance régulière provoque-t-elle avec 4\,096 ticks~?
\item Combien la mesure en a-t-elle compté, et pourquoi moins~?
\item Avec 256 ticks, la mesure en compte bien plus que la seule tendance n'en prévoit. Pourquoi~?
\item Que coûte un recentrage, d'après les mesures à 4\,096 et à 65\,536 ticks~?
\end{enumerate}

\medskip\noindent\textbf{Partie III --- Le compromis.}
\begin{enumerate}[resume]
\item Quelle largeur minimise le total mesuré de la journée, et que coûte-t-elle~?
\item Pourquoi le total remonte-t-il au-delà~?
\item Quelle mémoire cette largeur occupe-t-elle pour les deux côtés~?
\item Que dit le temps médian de mise à jour sur la largeur~?
\end{enumerate}

\medskip\noindent\textbf{Partie IV --- Le verdict.}
\begin{enumerate}[resume]
\item Énoncer le \emph{résultat nommé}~: la largeur de l'échelle face aux recentrages de la journée et à son coût total, et la largeur qui
le minimise.
\item Comment choisir la largeur pour chaque instrument~?
\item Que changerait un recentrage en arrière-plan, dans une seconde échelle~?
\item Pourquoi les tests doivent-ils inclure une échelle étroite~?
\item Comment un instrument à grand tick modifie-t-il la réponse~?
\item Que doit signaler le constructeur à la surveillance au sujet des recentrages~?
\item Que coûte la journée avec un arbre équilibré à la place~?
\item En une phrase~: qu'apporte une échelle dense, et contre quoi faut-il la protéger~?
\end{enumerate}
\end{problem}

\section{Questions d'entretien}

\begin{interviewq}[$\star$ \iqroles{développeur}]\label{iq:ll:the-order-book-builder:1}
Concevriez-vous un carnet d'ordres pour un flux market-by-order~? Quelles opérations doivent être rapides, et à quelle vitesse~?
\end{interviewq}

\begin{interviewq}[$\star\star$ \iqroles{développeur}]\label{iq:ll:the-order-book-builder:2}
Arbre, vecteur trié ou tableau de niveaux de prix~: quand choisiriez-vous chacun~?
\end{interviewq}

\begin{interviewq}[$\star\star$ \iqroles{développeur}]\label{iq:ll:the-order-book-builder:3}
Comment retrouveriez-vous des ordres par référence sans allouer de mémoire~? Comment supprimez-vous dans une table à \omterm{def:ll:the-order-book-builder:open}{adressage ouvert}~?
\end{interviewq}

\begin{interviewq}[$\star\star$ \iqroles{développeur}]\label{iq:ll:the-order-book-builder:4}
Votre carnet est croisé. Qu'a-t-il pu se passer, et que faites-vous~?
\end{interviewq}

\begin{interviewq}[$\star\star$ \iqroles{développeur}]\label{iq:ll:the-order-book-builder:5}
Comment testez-vous un \omterm{def:ll:the-order-book-builder:builder}{constructeur de carnet} pour lui faire confiance le jour qui compte~?
\end{interviewq}

\begin{interviewq}[$\star\star\star$ \iqroles{développeur, chercheur}]\label{iq:ll:the-order-book-builder:6}
Une échelle dense est rapide jusqu'à ce que le prix s'éloigne beaucoup. Quantifiez à partir de quand, et concevez en conséquence.
\end{interviewq}
